% =========================================================
% The Covering Relation and Hasse Diagrams
% =========================================================

\section{The Covering Relation}
\label{sec:covering-relation}

\begin{remark*}[Toolkit: The Covering Relation]
The covering relation identifies the irreducible steps of a partial order:
$y$ covers $x$ when $x < y$ and there is nothing strictly between them.
In a finite poset the entire order is determined by its covering relation,
which is encoded exactly by the Hasse diagram.
\end{remark*}

% ---------------------------------------------------------
% def:covering-relation
% ---------------------------------------------------------

\begin{definitionbox}{Definition (Covering Relation)}
\begin{definition}[Covering Relation]
\label{def:covering-relation}
Let $(P, \leq)$ be a partially ordered set and let $x, y \in P$ with
$x < y$.  The element $y$ \emph{covers} $x$, written $x \prec y$, if
there is no element of $P$ strictly between them:
\[
x \prec y
\;\coloneqq\;
x < y
\;\land\;
\neg\,\bigl(\exists z \in P,\; x < z < y\bigr).
\]
\end{definition}
\end{definitionbox}

\begin{remark*}[Standard quantified statement]
\[
x \prec y
\;\coloneqq\;
x < y
\;\land\;
\forall z \in P,\; x < z \leq y \Longrightarrow z = y.
\]
\end{remark*}

\begin{remark*}[Negated quantified statement]
\[
\neg(x \prec y)
\;\coloneqq\;
x \not< y \;\lor\; \exists z \in P,\; x < z < y.
\]
\end{remark*}

\begin{remark*}[Interpretation]
The covering relation captures the direct steps of the order --- the
comparisons that cannot be decomposed further.  In a finite poset every
strict comparison $x < y$ can be resolved into a finite chain of covers
$x = x_0 \prec x_1 \prec \cdots \prec x_n = y$, so the covering relation
alone determines the full order.  The Hasse diagram of a finite poset is
exactly the directed graph of the covering relation.  In a dense order
such as $\mathbb{Q}$ or $\mathbb{R}$, no element covers any other: between
any two elements there is always a third.
\end{remark*}

\begin{dependencies}
\hyperref[def:lattice-order-partial-order]{Partial order},
\hyperref[def:lattice-order-strict-partial-order]{strict partial order}.
\end{dependencies}

% ---------------------------------------------------------
% lem:order-iso-finite-characterizations
% ---------------------------------------------------------

\begin{lemmabox}{Lemma}
\begin{lemma}[Characterizations of Order-Isomorphisms for Finite Posets]
\label{lem:order-iso-finite-characterizations}
Let $(P, \leq_P)$ and $(Q, \leq_Q)$ be finite partially ordered sets and
let $\varphi : P \to Q$ be a bijection.  The following are equivalent:
\begin{enumerate}
  \item $\varphi$ is an order-isomorphism.
  \item For all $x, y \in P$:\; $x <_P y \;\Longleftrightarrow\; \varphi(x) <_Q \varphi(y)$.
  \item For all $x, y \in P$:\; $x \prec_P y \;\Longleftrightarrow\; \varphi(x) \prec_Q \varphi(y)$.
\end{enumerate}
\hyperref[prf:order-iso-finite-characterizations]{Go to proof.}
\end{lemma}
\end{lemmabox}

\begin{remark*}[Standard quantified statement]
\[
\operatorname{OrderIso}(\varphi, P, Q)
\;\Longleftrightarrow\;
\bigl(\forall x,y \in P,\; x <_P y \Leftrightarrow \varphi(x) <_Q \varphi(y)\bigr)
\;\Longleftrightarrow\;
\bigl(\forall x,y \in P,\; x \prec_P y \Leftrightarrow \varphi(x) \prec_Q \varphi(y)\bigr).
\]
\end{remark*}

\begin{remark*}[Interpretation]
In a finite poset the order relation, the strict order, and the covering
relation each determine the others completely.  This lemma shows that an
order-isomorphism can be detected at any one of these three levels.  The
covering characterisation is the most useful in practice: since Hasse
diagrams encode exactly the covering relation, two finite posets are
order-isomorphic precisely when their diagrams are isomorphic as directed
graphs.
\end{remark*}

\begin{dependencies}
\hyperref[def:order-isomorphism]{Order-isomorphism},
\hyperref[def:covering-relation]{covering relation}.
\end{dependencies}

% ---------------------------------------------------------
% prop:hasse-diagram-characterization
% ---------------------------------------------------------

\begin{propositionbox}{Proposition}
\begin{proposition}[Hasse Diagram Characterization for Finite Posets]
\label{prop:hasse-diagram-characterization}
Let $(P, \leq_P)$ and $(Q, \leq_Q)$ be finite partially ordered sets.
Then $P$ and $Q$ are order-isomorphic if and only if their Hasse diagrams
are isomorphic as directed graphs --- that is, identical up to relabeling
of elements.
\hyperref[prf:hasse-diagram-characterization]{Go to proof.}
\end{proposition}
\end{propositionbox}

\begin{remark*}[Standard quantified statement]
\[
P \cong Q
\;\Longleftrightarrow\;
\exists \varphi : P \xrightarrow{\;\sim\;} Q,\;
\forall x, y \in P,\;
x \prec_P y \Leftrightarrow \varphi(x) \prec_Q \varphi(y).
\]
\end{remark*}

\begin{remark*}[Interpretation]
The Hasse diagram is a complete invariant of the order structure of a
finite poset.  Checking order-isomorphism between finite posets reduces
to a graph isomorphism problem on the covering graphs --- two finite posets
with non-isomorphic Hasse diagrams cannot be order-isomorphic.
\end{remark*}

\begin{dependencies}
\hyperref[lem:order-iso-finite-characterizations]{Characterizations of order-isomorphisms for finite posets},
\hyperref[def:covering-relation]{covering relation}.
\end{dependencies}

